\section*{Chapitre \ref{ch:b2:comparison} --- Comparaison des fonctions}

\begin{solution}{exo:b2:comparison:1}
$\sqrt{x^2 + x + 1} = x\sqrt{1 + \tfrac1x + \tfrac{1}{x^2}} = x +
\frac12 + \frac38\cdot\frac1x + o\bigl(\frac1x\bigr)$ (développement
binomial~: $\frac12 u - \frac18 u^2$ avec $u = \frac1x +
\frac{1}{x^2}$ donne $\frac{1}{2x} + \frac{1}{2x^2} -
\frac{1}{8x^2} = \frac{1}{2x} + \frac{3}{8x^2}$, puis multiplier par
$x$).

$\ln(x^2 + x) - 2\ln x = \ln\bigl(1 + \tfrac1x\bigr) = \frac1x -
\frac{1}{2x^2} + o\bigl(\frac{1}{x^2}\bigr)$.

Troisième fonction~: développer chaque facteur,
\[
\frac{x + \sin x}{x - \ln x}
= \Bigl(1 + \frac{\sin x}{x}\Bigr)
\Bigl(1 + \frac{\ln x}{x} + \frac{(\ln x)^2}{x^2} +
O\Bigl(\frac{(\ln x)^3}{x^3}\Bigr)\Bigr).
\]
Ordonner les contributions sur l'échelle en $+\infty$~: $\frac{\ln x}{x}
\gg \frac{1}{x} \geq \bigl|\frac{\sin x}{x}\bigr| \gg \frac{(\ln
x)^2}{x^2}$. Les deux termes après le $1$ dominant sont donc
$\frac{\ln x}{x}$, puis le terme d'oscillation bornée $\frac{\sin
x}{x}$~:
\[
\frac{x + \sin x}{x - \ln x}
= 1 + \frac{\ln x}{x} + \frac{\sin x}{x}
+ O\Bigl(\frac{(\ln x)^2}{x^2}\Bigr).
\]
\end{solution}

\begin{solution}{exo:b2:comparison:2}
$f(t) = t^\alpha$ ($t \geq 1$).

$\alpha > -1$~: divergence, et d'après le~%
\cref{thm:b2:comparison:seriesintegral}~(2), $\sum_{k\leq n}
k^\alpha = \frac{n^{\alpha+1}}{\alpha+1} + C + o(1)$ si $\alpha <
0$ (où $f$ décroît)~; pour $\alpha \geq 0$ ($f$ croissante) le
même encadrement avec inégalités inversées donne $\sum_{k \leq n}
k^\alpha \sim \frac{n^{\alpha + 1}}{\alpha + 1}$.

$\alpha = -1$~: $H_n = \ln n + \gamma + o(1)$
(\cref{ex:b2:comparison:harmonic}).

$\alpha < -1$~: convergence, avec reste $\sum_{k > n} k^\alpha
\sim \frac{n^{\alpha+1}}{-(\alpha+1)}$ par l'encadrement (1) (les deux
bornes intégrales sont \omterm{def:b2:nvs:equivalent}{équivalentes} à cette valeur).
\end{solution}

\begin{solution}{exo:b2:comparison:3}
Soit $v_n = H_n - \ln n - \gamma \to 0$. Alors
\[
v_n - v_{n+1} = \ln\frac{n+1}{n} - \frac{1}{n+1}
= \Bigl(\frac1n - \frac{1}{2n^2}\Bigr) - \Bigl(\frac1n -
\frac{1}{n^2}\Bigr) + O\Bigl(\frac{1}{n^3}\Bigr)
= \frac{1}{2n^2} + O\Bigl(\frac{1}{n^3}\Bigr),
\]
en utilisant $\frac{1}{n+1} = \frac1n - \frac{1}{n^2} + O(n^{-3})$.
Comme $v_n \to 0$, en télescopant le reste~:
\[
v_n = \sum_{k \geq n} (v_k - v_{k+1})
= \sum_{k\geq n} \Bigl(\frac{1}{2k^2} + O(k^{-3})\Bigr)
= \frac{1}{2n} + O\Bigl(\frac{1}{n^2}\Bigr),
\]
d'après le~\cref{thm:b2:comparison:seriesintegral}~(1) appliqué à
$t^{-2}$ (reste $\sim \frac1n$, divisé par deux) et à $t^{-3}$. D'où
$H_n = \ln n + \gamma + \frac{1}{2n} + o(\frac1n)$.
\end{solution}

\begin{solution}{exo:b2:comparison:4}
Stirling trois fois~:
\[
\frac{(3n)!}{(n!)^3}
\sim \frac{\sqrt{6\pi n}\,(3n/\eu)^{3n}}
{(2\pi n)^{3/2}\,(n/\eu)^{3n}}
= \frac{\sqrt{6}\; 27^{\,n}}{2\pi n} \cdot
\frac{1}{\sqrt{2\pi n}}\cdot\sqrt{2\pi n}\;
= \frac{\sqrt3\,27^n}{2\pi n} .
\]
(Attentivement~: $\frac{\sqrt{6\pi n}}{(2\pi n)^{3/2}} =
\frac{\sqrt6}{(2\pi n)\sqrt{2\pi n}}\sqrt{\pi n} =
\frac{\sqrt3}{2\pi n}$.)

$\dfrac{n!}{n^n} \sim \sqrt{2\pi n}\,\eu^{-n}$.

$\sqrt[n]{n!} = \exp\bigl(\frac{\ln n!}{n}\bigr)$ avec $\ln n! = n\ln
n - n + \frac12\ln(2\pi n) + o(1)$~:
\[
\sqrt[n]{n!} = \exp\Bigl(\ln n - 1 + \frac{\ln(2\pi n)}{2n} +
o\Bigl(\frac{\ln n}{n}\Bigr)\Bigr)
= \frac{n}{\eu}\Bigl(1 + \frac{\ln(2\pi n)}{2n} +
o\Bigl(\frac{\ln n}{n}\Bigr)\Bigr).
\]
\end{solution}

\begin{solution}{exo:b2:comparison:5}
$g(x) = x^n + x - 1$ croît strictement sur $\intcc{0}{1}$ de $-1$
à $1$~: racine unique $x_n$. Comme $x_n^n = 1 - x_n \in
\intoo{0}{1}$~: si $x_n \leq c < 1$ le long d'une sous-suite, alors
$x_n^n \leq c^n \to 0$, donc $1 - x_n \to 0$~: contradiction avec $x_n
\leq c$. D'où $x_n \to 1$.

Écrivons $x_n = 1 - \varepsilon_n$, $\varepsilon_n \to 0^+$.
L'équation s'écrit $(1 - \varepsilon_n)^n = \varepsilon_n$,
c'est-à-dire
\[
n\ln(1 - \varepsilon_n) = \ln \varepsilon_n
\quad\Longrightarrow\quad
-n\varepsilon_n\bigl(1 + o(1)\bigr) = \ln\varepsilon_n .
\]
Donc $n\varepsilon_n = -\ln\varepsilon_n\,(1 + o(1)) \to +\infty$, et
en prenant de nouveau les logarithmes~: $\ln n + \ln\varepsilon_n =
\ln(-\ln\varepsilon_n) + o(1)$. Comme $\ln(-\ln \varepsilon_n) =
o(\ln(1/\varepsilon_n))$, cela donne $\ln\varepsilon_n \sim -\ln
n$, d'où $\varepsilon_n = \frac{-\ln\varepsilon_n}{n}(1 + o(1))
\sim \frac{\ln n}{n}$~:
\[
x_n = 1 - \frac{\ln n}{n} + o\Bigl(\frac{\ln n}{n}\Bigr) .
\]
\end{solution}

\begin{solution}{exo:b2:comparison:6}
Relation exacte~: $\cot y_n = x_n = n\pi + \frac\pi2 - y_n$, avec $y_n
\sim \frac{1}{n\pi}$ (\cref{ex:b2:comparison:tan}). Développer $\cot y
= \frac1y - \frac y3 + O(y^3)$~:
\[
\frac{1}{y_n} - \frac{y_n}{3} + O(y_n^3) = n\pi + \frac\pi2 - y_n
\quad\Longrightarrow\quad
\frac{1}{y_n} = n\pi + \frac\pi2 + O\Bigl(\frac1n\Bigr),
\]
(les termes $-y_n$ et $-\frac{y_n}{3}$ sont $O(\frac1n)$). Inverser~:
\[
y_n = \frac{1}{n\pi}\cdot\frac{1}{1 + \frac{1}{2n} + O(n^{-2})}
= \frac{1}{n\pi}\Bigl(1 - \frac{1}{2n} + O\Bigl(\frac{1}{n^2}\Bigr)\Bigr)
= \frac{1}{n\pi} - \frac{1}{2n^2\pi} + O\Bigl(\frac{1}{n^3}\Bigr).
\]
D'où
\[
x_n = n\pi + \frac{\pi}{2} - y_n
= n\pi + \frac\pi2 - \frac{1}{n\pi} + \frac{1}{2n^2\pi} +
o\Bigl(\frac{1}{n^2}\Bigr).
\]
\end{solution}

\begin{solution}{exo:b2:comparison:7}
Séparons les deux plus grands termes~:
\[
\sum_{k=0}^{n} k! = n! + (n-1)! + \sum_{k \leq n-2} k! ,
\qquad
\sum_{k\leq n-2} k! \leq (n-1)\,(n-2)! = (n-1)! .
\]
Donc $1 \leq \frac{1}{n!}\sum k! \leq 1 + \frac{2}{n}$~: la limite est
$1$. En raffinant~: $\frac{(n-1)!}{n!} = \frac1n$ et la majoration
grossière $\sum_{k \leq n-2}k! \leq (n-1)!$ peut être affinée de la
même manière~: $\sum_{k\leq n-2} k! = (n-2)!\,(1 + O(\frac1n)) =
O\bigl(\frac{n!}{n^2}\bigr)$. D'où
\[
\sum_{k=0}^{n} k! = n!\Bigl(1 + \frac1n + O\Bigl(\frac{1}{n^2}\Bigr)\Bigr).
\]
\end{solution}

\begin{solution}{exo:b2:comparison:8}
$(u_n)$ croît~; si elle était bornée elle convergerait vers $\ell$
avec $\ell = \ell + \frac1\ell$~: absurde. Donc $u_n \to \infty$.

Carrés~: $u_{n+1}^2 = u_n^2 + 2 + u_n^{-2}$, donc
\[
u_n^2 = u_0^2 + 2n + \sum_{k=0}^{n-1} \frac{1}{u_k^2} .
\]
La somme est $o(n)$ (les termes tendent vers $0$, Cesàro), donc $u_n^2
\sim 2n$ et $u_n \sim \sqrt{2n}$.

Raffinement~: $\frac{1}{u_k^2} \sim \frac{1}{2k}$, donc par comparaison
(\cref{thm:b2:comparison:seriesintegral}, ou équivalents de sommes
partielles de séries positives) $\sum_{k<n} u_k^{-2} \sim \frac12 \ln
n$. D'où
\[
u_n^2 = 2n + \frac{\ln n}{2}\,(1 + o(1)) + O(1)
\quad\Longrightarrow\quad
u_n = \sqrt{2n}\sqrt{1 + \frac{\ln n}{4n} + o\Bigl(\frac{\ln
n}{n}\Bigr)}
= \sqrt{2n}\Bigl(1 + \frac{\ln n}{8n} + o\Bigl(\frac{\ln
n}{n}\Bigr)\Bigr).
\]
\end{solution}

\begin{solution}{exo:b2:comparison:9}
Factoriser $n$ et poser $t = \ln n$~:
\[
S_n = \frac1n \sum_{k=1}^{n} \frac{1}{1 + t\,\frac kn} .
\]
Pour $t$ fixé, la somme est une somme de Riemann de $u \mapsto \frac{1}{1 +
tu}$ sur $\intcc{0}{1}$~; la fonction est monotone en $u$, donc la
somme de Riemann est encadrée par l'intégrale décalée d'une maille~:
\[
\int_0^1 \frac{\dd u}{1 + tu} - \frac1n
\leq S_n \leq \int_0^1 \frac{\dd u}{1 + tu} + \frac1n
\]
(comparaison des sommes de Riemann d'une fonction monotone avec son
intégrale, valable pour chaque $n$ avec son propre $t = \ln n$).
Maintenant $\int_0^1 \frac{\dd u}{1 + tu} = \frac{\ln(1 + t)}{t}$, et
$\frac1n = o\bigl(\frac{\ln t}{t}\bigr)$~: d'où
\[
S_n = \frac{\ln(1 + \ln n)}{\ln n} + O\Bigl(\frac 1n\Bigr)
\;\sim\; \frac{\ln\ln n}{\ln n} .
\]
\end{solution}

\begin{solution}{exo:b2:comparison:10}
Identité~: $(\ln n)^{\ln n} = \eu^{\ln n\,\ln\ln n} =
\bigl(\eu^{\ln n}\bigr)^{\ln\ln n} = n^{\ln\ln n}$. Classement~:
comparer les logarithmes. $\ln(n^2) = 2\ln n$~; $\ln\bigl((\ln
n)^{\ln n}\bigr) = \ln n\ln\ln n$~; $\ln(2^n) = n\ln2$~;
$\ln(n!) = n\ln n - n + O(\ln n)$ (Stirling, ou l'encadrement plus
grossier $\ln n! \sim n\ln n$)~; $\ln(n^n) = n\ln n$. Comme
$2\ln n = o(\ln n\ln\ln n)$, $\ln n\ln\ln n = o(n)$, $n\ln 2 =
o(n\ln n - n)$, et $n \ln n - n \sim n\ln n$ mais $n! / n^n \to
0$ (la différence des logarithmes est $-n + O(\ln n) \to -\infty$)~:
\[
n^2 = o\bigl((\ln n)^{\ln n}\bigr),\quad
(\ln n)^{\ln n} = o(2^n),\quad
2^n = o(n!),\quad
n! = o(n^n).
\]
(Pour chaque étape~: la différence des logarithmes tend vers
$+\infty$, donc le rapport tend vers $0$.)
\end{solution}

\begin{solution}{exo:b2:comparison:11}
Décomposer $\frac1{k^2} = \frac1{k(k+1)} + \frac1{k^2(k+1)}$ et
sommer pour $k > n$~:
\[
\sum_{k>n}\frac1{k^2} = \frac1{n+1} +
\sum_{k>n}\frac{1}{k^2(k+1)} ,
\]
la première somme se télescopant exactement ($\frac1{k(k+1)} = \frac1k -
\frac1{k+1}$). Pour la seconde~: $\frac{1}{k^2(k+1)} = \frac1{k^3}
+ O\bigl(\frac1{k^4}\bigr)$ (car $\frac{1}{k^2(k+1)} -
\frac1{k^3} = \frac{-1}{k^3(k+1)}$), et par la
comparaison intégrale $\sum_{k>n}\frac1{k^3} = \frac1{2n^2} +
O\bigl(\frac1{n^3}\bigr)$, $\sum_{k>n}\frac1{k^4} =
O\bigl(\frac1{n^3}\bigr)$. D'où
\[
\sum_{k>n}\frac1{k^2}
= \frac1{n+1} + \frac{1}{2n^2} + O\Bigl(\frac1{n^3}\Bigr)
= \frac1n - \frac1{n^2} + \frac{1}{2n^2} +
O\Bigl(\frac1{n^3}\Bigr)
= \frac1n - \frac{1}{2n^2} + O\Bigl(\frac1{n^3}\Bigr),
\]
en utilisant $\frac1{n+1} = \frac1n - \frac1{n^2} +
O\bigl(\frac1{n^3}\bigr)$.
\end{solution}

\begin{solution}{exo:b2:comparison:12}
$(u_n)$ croît~; si elle était bornée elle convergerait vers un
$\ell$ fini avec $\ell = \ell + \eu^{-\ell}$~: impossible. Donc
$u_n \to \infty$. Soit $v_n = \eu^{u_n} \to \infty$~: alors
\[
v_{n+1} = \eu^{u_n + \eu^{-u_n}} = v_n\,\eu^{1/v_n}
= v_n\Bigl(1 + \frac1{v_n} + \frac1{2v_n^2} +
O\bigl(v_n^{-3}\bigr)\Bigr)
= v_n + 1 + \frac{1}{2v_n} + O\bigl(v_n^{-2}\bigr).
\]
En sommant $v_{k+1} - v_k = 1 + O(1)$ on obtient d'abord $v_n = n +
O(n)$, d'où $v_n \geq cn$ à partir d'un certain rang~; en re-sommant
avec $\frac1{2v_k} = O(\frac1k)$ on obtient $v_n = n + O(\ln n)$. Une
passe de plus~: $\frac{1}{2v_k} = \frac{1}{2k}\bigl(1 +
O\bigl(\tfrac{\ln k}k\bigr)\bigr)$, donc
\[
v_n = n + \sum_{k<n}\frac1{2k} + O(1) = n + \frac{\ln n}2 +
O(1).
\]
Enfin $u_n = \ln v_n = \ln n + \ln\Bigl(1 + \frac{\ln n}{2n} +
O\bigl(\tfrac1n\bigr)\Bigr) = \ln n + \frac{\ln n}{2n} +
O\bigl(\tfrac1n\bigr)$.
\end{solution}

\begin{solution}{pb:b2:comparison:1}
\textbf{1.} En soustrayant les deux développements~: $\sum_i (c_i -
c_i')\varphi_i = o(\varphi_k)$. Si un coefficient diffère, soit
$i_0$ le premier~: en divisant par $\varphi_{i_0}$ et en utilisant
$\varphi_j = o(\varphi_{i_0})$ pour $j > i_0$ on obtient $c_{i_0} -
c_{i_0}' = o(1)$~: nul, contradiction. Pour le développement~: avec
$u = \frac{\ln x}x \to 0$,
\[
\frac{1}{x - \ln x} = \frac1x\cdot\frac{1}{1 - u}
= \frac1x\bigl(1 + u + u^2 + O(u^3)\bigr)
= \frac1x + \frac{\ln x}{x^2} + \frac{(\ln x)^2}{x^3} +
o\Bigl(\frac{(\ln x)^2}{x^3}\Bigr).
\]
Aucun terme $\frac c{x^2}$ n'apparaît parce que le développement est
une série géométrique en $u = \frac{\ln x}{x}$~: chaque terme porte
autant de puissances de $\ln x$ que de $\frac1x$ au-delà de la
première~; l'échelon d'échelle $\frac1{x^2}$ (coefficient de $(\ln
x)^0$) est simplement absent, de coefficient $0$.

\textbf{2.} $f(x) = \eu^x + x$ est \omterm{def:b2:metric:continuity}{continue}, strictement
croissante, de limites $-\infty$ et $+\infty$~: une bijection
$\R \to \R$, donc $x_n = f^{-1}(n)$ existe et est unique, et
$x_n \to +\infty$ ($f^{-1}$ croît vers $+\infty$). De
$\eu^{x_n} = n - x_n$~: $x_n = \ln(n - x_n) \leq \ln n$, donc
$x_n/n \to 0$ et $x_n = \ln n + \ln(1 - x_n/n) = \ln n + o(1)
\sim \ln n$.

\textbf{3.} Posons $u_n = x_n/n$. Deuxième passe~: $u_n = \frac{\ln n
+ o(1)}{n}$, donc
\[
x_n = \ln n + \ln(1 - u_n) = \ln n - u_n + O(u_n^2)
= \ln n - \frac{\ln n}{n} + o\Bigl(\frac{\ln n}n\Bigr).
\]
Troisième passe~: maintenant $u_n = \frac{\ln n}{n} - \frac{\ln n}{n^2} +
o\bigl(\frac{\ln n}{n^2}\bigr)$, et $\ln(1 - u_n) = -u_n -
\frac{u_n^2}2 + O(u_n^3)$~:
\[
x_n = \ln n - \frac{\ln n}n + \frac{\ln n}{n^2}
- \frac{(\ln n)^2}{2n^2} + o\Bigl(\frac{(\ln
n)^2}{n^2}\Bigr)
= \ln n - \frac{\ln n}{n} - \frac{(\ln n)^2}{2n^2} +
o\Bigl(\frac{(\ln n)^2}{n^2}\Bigr),
\]
le terme $\frac{\ln n}{n^2}$ étant absorbé dans
$o\bigl(\frac{(\ln n)^2}{n^2}\bigr)$.

\textbf{4.} En $n = 1000$~: $\ln 1000 \approx 6.90776$ (erreur
$7\cdot10^{-3}$)~; deux termes~: $6.90085$ (erreur
$2\cdot10^{-5}$)~; trois termes~: $6.90082$ (erreur inférieure à
$10^{-5}$), contre $x_{1000} \approx 6.90083$. Chaque passe achète
grossièrement le facteur prévu $\frac{\ln n}{n}$.

\textbf{5.} Deux intégrations par parties, en partant de la droite~:
avec $\frac{\dd}{\dd t}\bigl[\tfrac12t(1-t)\bigr] = \tfrac12 -
t$ et $t(1-t)$ s'annulant aux deux extrémités,
\[
\frac12\int_0^1 t(1-t)g''(t)\dd t
= -\int_0^1\Bigl(\frac12 - t\Bigr)g'(t)\dd t
= -\Bigl[\Bigl(\frac12 - t\Bigr)g\Bigr]_0^1 - \int_0^1 g
= \frac{g(0) + g(1)}2 - \int_0^1 g .
\]
Réarrangé, c'est l'identité annoncée.

\textbf{6.} Calculer l'accroissement, puis appliquer la question 5 à
$g(t) = f(n + t)$~:
\begin{align*}
E_{n+1} - E_n
&= f(n{+}1) - \int_n^{n+1}\!f - \frac{f(n{+}1) - f(n)}2 \\
&= \frac{f(n) + f(n{+}1)}2 - \int_n^{n+1}\!f
= \frac12\int_0^1 t(1-t)f''(n+t)\dd t .
\end{align*}
Comme $0 \leq t(1-t) \leq \frac14$~: $\abs{E_{n+1} - E_n} \leq
\frac18\int_n^{n+1}\abs{f''}$, dont la somme sur $n$ converge par
hypothèse~: $(E_n)$ converge (accroissements absolument sommables)
vers un $E$, avec
\[
\abs{E - E_n} \leq \sum_{k\geq n}\abs{E_{k+1} - E_k} \leq
\frac18\int_n^\infty\abs{f''} .
\]

\textbf{7.} $f(t) = \frac1t$~: $f''(t) = \frac2{t^3}$,
$\int_1^\infty\abs{f''} = 1 < \infty$. Question 6~:
\[
H_n = \ln n + \frac{1 + \frac1n}{2} + E + (E_n - E)
= \ln n + \Bigl(E + \frac12\Bigr) + \frac1{2n} +
\varepsilon_n,
\]
avec $\abs{\varepsilon_n} = \abs{E_n - E} \leq
\frac18\int_n^\infty\frac{2\dd t}{t^3} = \frac1{8n^2}$.
En comparant avec $H_n = \ln n + \gamma + o(1)$
(\cref{ex:b2:comparison:harmonic}) on identifie $E + \frac12 =
\gamma$.

\textbf{8.} À partir de la formule d'accroissement de la question 6,
\[
\varepsilon_n = E_n - E = -\sum_{k\geq n}\frac12\int_0^1
t(1-t)\,\frac{2\,\dd t}{(k+t)^3}
= -\sum_{k \geq n}\Bigl(\frac1{k^3}\int_0^1t(1-t)\dd t +
O\Bigl(\frac1{k^4}\Bigr)\Bigr),
\]
en utilisant $\frac{1}{(k+t)^3} = \frac1{k^3} +
O\bigl(\frac1{k^4}\bigr)$ uniformément pour $t \in \intcc01$. Avec
$\int_0^1 t(1-t) = \frac16$ et
$\sum_{k\geq n}\frac1{k^3} \sim \frac{1}{2n^2}$
(\cref{thm:b2:comparison:seriesintegral})~:
\[
\varepsilon_n = -\frac16\cdot\frac{1}{2n^2} +
o\Bigl(\frac1{n^2}\Bigr) = -\frac{1}{12n^2} +
o\Bigl(\frac{1}{n^2}\Bigr).
\]

\textbf{9.} $f = \ln$~: $f''(t) = -\frac1{t^2}$, absolument
intégrable. La question 6 donne
\[
\ln n! = \int_1^n\ln t\,\dd t + \frac{\ln n}2 + E + O\Bigl(
\frac1{8}\int_n^\infty\frac{\dd t}{t^2}\Bigr)
= \Bigl(n + \frac12\Bigr)\ln n - n + 1 + E +
O\Bigl(\frac1n\Bigr),
\]
donc $d_n = 1 + E + O\bigl(\frac1n\bigr)$~: convergence de $(d_n)$
--- étape 1 du~\cref{thm:b2:comparison:stirling} --- plus le
taux $O(1/n)$. (La valeur de la limite donnée par Stirling fournit $E =
\ln\sqrt{2\pi} - 1$.)

\textbf{10.} $f(t) = t^{-1/2}$~: $f''(t) = \frac34 t^{-5/2}$,
absolument intégrable. Question 6~:
\[
\sum_{k=1}^n \frac1{\sqrt k}
= 2\sqrt n - 2 + \frac{1 + \frac1{\sqrt n}}2 + E +
O\bigl(n^{-3/2}\bigr)
= 2\sqrt n + c + \frac{1}{2\sqrt n} +
O\bigl(n^{-3/2}\bigr),
\]
avec $c = E - \frac32$. En $n = 10^4$~: $2\sqrt n = 200$, $c
\approx -1.46035$, $\frac1{2\sqrt n} = 0.005$~: valeur prédite
$198.54465$, et en effet $\sum_{k\leq10^4}k^{-1/2} =
198.544645\dots$ --- trois termes, sept chiffres.

\textbf{11.} $t \mapsto t\ln t$ est \omterm{def:b2:metric:continuity}{continue} et strictement
croissante sur $\intco1\infty$ (dérivée $\ln t + 1 \geq 1$),
de $0$ à $+\infty$~: un unique $x_n$ existe, et $x_n \to
\infty$ (sinon $x_n\ln x_n$ resterait borné). En prenant les
logarithmes dans $x_n\ln x_n = n$~: $\ln x_n + \ln\ln x_n = \ln n$~;
comme $\ln\ln x_n = o(\ln x_n)$, en divisant par $\ln x_n$ on obtient
$\frac{\ln n}{\ln x_n} \to 1$~: $\ln x_n \sim \ln n$.

\textbf{12.} De $x_n = \frac{n}{\ln x_n}$ et $\ln x_n \sim
\ln n$~: $x_n \sim \frac{n}{\ln n}$. Passe suivante~: $\ln\ln x_n =
\ln\bigl(\ln n\,(1 + o(1))\bigr) = \ln\ln n + o(1)$, donc $\ln
x_n = \ln n - \ln\ln n + o(1)$ et
\[
x_n = \frac{n}{\ln n - \ln\ln n + o(1)}
= \frac{n}{\ln n}\cdot\frac{1}{1 - \frac{\ln\ln n +
o(1)}{\ln n}}
= \frac{n}{\ln n}\Bigl(1 + \frac{\ln\ln n}{\ln n} +
o\Bigl(\frac{\ln\ln n}{\ln n}\Bigr)\Bigr).
\]

\textbf{13.} En $n = 10^6$~: $\frac{n}{\ln n} \approx 72\,382$
(à $18\%$ près), deux termes donnent $\approx 86\,140$ (à
$1.9\%$ près), contre le vrai $x \approx 87\,848$. Le gain par
passe n'est que le facteur $\frac{\ln\ln n}{\ln n} \approx
\frac{2.63}{13.8} \approx 0.19$~: les échelles logarithmiques
convergent avec une lenteur exaspérante --- un fait de la vie partout
où interviennent les nombres premiers.

\textbf{14.} Il y a exactement $n$ nombres premiers $\leq p_n$ (à
savoir $p_1, \dots, p_n$)~: $\pi(p_n) = n$. Le théorème des nombres
premiers (admis~; volume de troisième année) donne $n = \pi(p_n) \sim
\frac{p_n}{\ln p_n}$, c'est-à-dire $p_n \sim n\ln p_n$~: c'est
l'équation $x\ln x \approx n$ lue à l'envers. En prenant les
logarithmes~: $\ln p_n = \ln n + \ln\ln p_n + o(1)$, et $\ln\ln p_n =
o(\ln p_n)$ force $\ln p_n \sim \ln n$ comme à la question 11.
En substituant~:
\[
p_n \sim n\ln p_n = n\,\ln n\,\frac{\ln p_n}{\ln n} \sim n\ln
n .
\]

\textbf{15.} (a) Fixons $\varepsilon > 0$~; pour $k$ grand, $(1 -
\varepsilon)k\ln k \leq p_k \leq (1 + \varepsilon)k\ln k$. Par
comparaison avec le $t\ln t$ croissant
(encadrement de type \cref{thm:b2:comparison:seriesintegral}),
$\sum_{k\leq n}k\ln k = \int_1^n t\ln t\,\dd t + O(n\ln n) =
\frac{n^2\ln n}2 - \frac{n^2}4 + O(n\ln n) \sim
\frac{n^2\ln n}2$. D'où $\sum_{k\leq n}p_k = \frac{n^2\ln
n}{2}(1 + O(\varepsilon) + o(1))$ pour tout $\varepsilon$~:
$\sum_{k\leq n}p_k \sim \frac{n^2\ln n}2$. (b) Par le théorème des
nombres premiers, parmi les entiers jusqu'à $10^{100}$ une
proportion $\sim \frac{1}{\ln 10^{100}} = \frac1{230.26\dots}$
sont premiers~: un entier à $100$ chiffres tiré uniformément au hasard
est premier avec probabilité environ $\frac1{230}$.

\textbf{16.} Sur $\intoo{n\pi}{n\pi + \frac\pi2}$, $g(x) = \tan
x - \frac1x$ est \omterm{def:b2:metric:continuity}{continue} et strictement croissante ($g' = 1 +
\tan^2x + \frac1{x^2} > 0$), avec $g \to -\frac1{n\pi} < 0$ à
l'extrémité gauche et $g \to +\infty$ à droite~: exactement une
racine $x_n$. Comme $\tan z_n = \tan x_n = \frac1{x_n} \to 0$
avec $z_n \in \intoo{0}{\frac\pi2}$~: $z_n \to 0^+$.

\textbf{17.} $\tan z_n \sim z_n$ et $\frac1{x_n} \sim
\frac1{n\pi}$~: $z_n \sim \frac1{n\pi}$.

\textbf{18.} $z_n = \arctan\frac1{x_n}$ et $\arctan u = u +
O(u^3)$. Avec $z_n = O(\frac1n)$~:
\[
\frac1{x_n} = \frac{1}{n\pi}\cdot\frac1{1 + \frac{z_n}{n\pi}}
= \frac1{n\pi} - \frac{z_n}{n^2\pi^2} +
O\Bigl(\frac1{n^4}\Bigr)
= \frac1{n\pi} + O\Bigl(\frac1{n^3}\Bigr),
\]
donc $z_n = \frac1{n\pi} + O\bigl(\frac1{n^3}\bigr)$~: l'échelon
$\frac{c}{n^2}$ porte le coefficient $0$, parce que la
première correction à $\frac1{x_n}$ est elle-même de taille
$\frac{z_n}{n^2} = O(n^{-3})$.

\textbf{19.} Insérer $z_n = \frac1{n\pi} + O(n^{-3})$ dans
l'affichage précédent~:
\[
\frac{1}{x_n} = \frac{1}{n\pi} - \frac{1}{n^3\pi^3} +
O\Bigl(\frac1{n^5}\Bigr),
\]
puis $z_n = \arctan\frac1{x_n} = \frac1{x_n} -
\frac{1}{3}\Bigl(\frac1{x_n}\Bigr)^3 + O\Bigl(\frac1{n^5}\Bigr)
= \frac1{n\pi} - \frac{1}{n^3\pi^3} - \frac{1}{3n^3\pi^3} +
O\Bigl(\frac1{n^5}\Bigr)$~:
\[
x_n = n\pi + \frac{1}{n\pi} - \frac{4}{3\pi^3n^3} +
O\Bigl(\frac1{n^5}\Bigr).
\]

\textbf{20.} En $n = 3$~: un terme $9.53088$, trois termes
$9.52929$, vraie racine $9.52933$~: erreurs $1.5\cdot10^{-3}$ et
$5\cdot10^{-5}$. Contraste~: pour $\tan x = x$ la racine doit rendre
$\tan$ énorme, donc elle se colle à l'extrémité \emph{droite} $n\pi +
\frac\pi2$ de la fenêtre, à distance $\sim\frac1{n\pi}$
avant l'asymptote~; pour $x\tan x = 1$ la racine doit rendre
$\tan$ minuscule, donc elle se place juste après l'extrémité
\emph{gauche} $n\pi$, à distance $\sim\frac1{n\pi}$ après le zéro.
Même méthode, géographie en miroir.

\textbf{21.} $\sin u < u$ sur $\intoo0\pi$ et $\sin$ envoie
$\intoo0\pi$ dans $\intoc01 \subseteq \intoo0\pi$~: après une
étape $u_1 \in \intoc{0}{1}$, puis $(u_n)$ décroît et est
minorée par $0$~: elle converge vers un point fixe de $\sin$,
c'est-à-dire vers $0$. Développement~: $\sin u = u(1 - \frac{u^2}6 +
o(u^2))$, donc
\[
\frac{1}{u_{n+1}^2} - \frac1{u_n^2}
= \frac{1}{u_n^2}\Bigl(\bigl(1 - \tfrac{u_n^2}6 +
o(u_n^2)\bigr)^{-2} - 1\Bigr)
= \frac{1}{u_n^2}\Bigl(\frac{u_n^2}{3} + o(u_n^2)\Bigr)
\longrightarrow \frac13 .
\]

\textbf{22.} Par Cesàro (volume de première année), la moyenne des
accroissements converge vers la même limite~:
\[
\frac{1}{n}\cdot\frac{1}{u_n^2}
= \frac1n\Bigl(\frac1{u_0^2} + \sum_{k=0}^{n-1}
\Bigl(\frac1{u_{k+1}^2} - \frac1{u_k^2}\Bigr)\Bigr)
\longrightarrow \frac13 ,
\]
donc $u_n^2 \sim \frac3n$ et, tous les termes étant positifs, $u_n
\sim \sqrt{3/n}$.

\textbf{23.} Par la question 7, $\abs{H_n - \ln n - \gamma -
\frac1{2n}} \leq \frac1{8n^2}$. En $n = 10^6$ cette borne vaut
$\frac{1}{8\cdot10^{12}} = 1.25\cdot10^{-13}$~: trois termes calculés
livrent la somme harmonique d'un million de termes à treize
chiffres, avec un certificat d'erreur pleinement rigoureux --- tout
l'intérêt d'une formule asymptotique à reste explicite.

\textbf{24.} (a) Vrai~: $\ln u_n - \ln v_n = \ln\frac{u_n}{v_n}
\to 0$ tandis que $\ln v_n \to +\infty$, donc le rapport des
logarithmes tend vers $1$. (b) Faux~: $u_n = n + 1 \sim v_n = n$, mais
$\eu^{u_n}/\eu^{v_n} = \eu \neq 1$. L'équivalence tolère des
erreurs additives $o(1)$ dans l'exposant, non $O(1)$. (c) Faux~:
$f(x) = x + \sin(x^2) \sim g(x) = x$ en $+\infty$, mais $f'(x)
= 1 + 2x\cos(x^2)$ oscille sans borne tandis que $g' = 1$~:
les dérivées de fonctions \omterm{def:b2:nvs:equivalent}{équivalentes} n'ont aucune raison d'être
comparables.

\textbf{25.} La boucle du~\cref{met:b2:comparison:implicit} a tourné
identiquement trois fois~: localiser la racine, extraire un terme
grossier, le réinjecter pour l'\omterm{def:b2:structures:generated}{ordre} suivant --- sur $\eu^x + x = n$
(partie I), sur $x\ln x = n$ (partie III), sur $x\tan x = 1$ (partie
IV). La correction trapèze fait passer la comparaison série--intégrale
de «~la différence converge~» à un terme explicite $\frac{f(1) +
f(n)}2$ avec un reste certifié $O(\int_n^\infty
\abs{f''})$ --- des constantes et des barres d'erreur au lieu de la
seule convergence. Le pont vers les nombres premiers est de l'inversion
pure~: le théorème des nombres premiers dit $\pi(x)\ln x \approx x$,
donc $p_n$, défini par $\pi(p_n) = n$, résout une équation $x\ln x = n$
--- et en hérite l'asymptotique. La règle (a) de la question 24 a
légitimé chaque passage de $u_n \sim v_n$ à $\ln u_n \sim
\ln v_n$ (questions 11, 14)~; la fausseté de (b) est la raison pour
laquelle on n'exponentie jamais les équivalences. Sommets~: la formule
d'Euler--Maclaurin au premier \omterm{def:b2:structures:generated}{ordre} (question 6), et la loi asymptotique
$p_n \sim n\ln n$ du $n$-ième nombre premier.
\end{solution}
